\begin{thebibliography}{99}

\bibitem{ayto_albacete_2021} 
Ayuntamiento de Albacete (2021) 
\textit{Guía de buenas prácticas de accesibilidad universal del Ayuntamiento de Albacete}. 
Albacete: ACCESIA y COCEMFE Albacete.

\bibitem{ayto_madrid_movilidad_2018} 
Ayuntamiento de Madrid (2018) 
«Ordenanza de Movilidad Sostenible, de 5 de octubre de 2018», 
\textit{Boletín Oficial del Ayuntamiento de Madrid}, nº 8259.

\bibitem{ayto_madrid_limpieza_2022} 
Ayuntamiento de Madrid (2022) 
«Ordenanza de Limpieza de los Espacios Públicos y Gestión de Residuos, de 26 de mayo de 2022», 
\textit{Boletín Oficial del Ayuntamiento de Madrid}, nº 9152.

\bibitem{cscae_fundacion_once_2024} 
CSCAE y Fundación ONCE (2024) 
\textit{Documenta 1.2 - Ciudad y Territorio justo: Accesibilidad universal}. 
Madrid: Consejo Superior de los Colegios de Arquitectos de España y Fundación ONCE. 
Disponible en: \url{https://observatorio2030.com/sites/default/files/2025-01/Documenta\%201.2\%20-\%20Ciudad\%20y\%20Territorio\%20justo\%20-\%20Accesibilidad\%20universal\%20\%28Informe\%20GT1.2\%29\%20\%5BAccesible\%5D_1.pdf}.

\bibitem{dacunha_2015} 
da Cunha, I. (2015) 
\textit{El trabajo de fin de grado y de máster: Redacción, defensa y publicación}. 
Barcelona: Editorial UOC.

\bibitem{fundacion_once_2022} 
Fundación ONCE (2022) 
\textit{Resumen ejecutivo: Índice de accesibilidad e inclusión. Ciudades de España 2022}. 
Madrid: Fundación ONCE. 
Disponible en: \url{https://biblioteca.fundaciononce.es/system/files/resumen_ejecutivo_indice_de_accesibilidad_e_inclusion._ciudades_de_espana_2022_def.pdf}.

\bibitem{ine_edad_2020} 
Instituto Nacional de Estadística (INE) (2022) 
\textit{Encuesta sobre Discapacidad, Autonomía Personal y Situaciones de Dependencia (EDAD 2020)}. 
Nota de prensa. Madrid: INE. 
Disponible en: \url{https://www.ine.es/prensa/edad_2020_p.pdf}.

\bibitem{rdl_1_2013} 
Jefatura del Estado (2013) 
«Real Decreto Legislativo 1/2013, de 29 de noviembre, por el que se aprueba el Texto Refundido de la Ley General de derechos de las personas con discapacidad y de su inclusión social», 
\textit{Boletín Oficial del Estado}, nº 289, págs. 95635--95673.

\bibitem{ministerio_derechos_sociales_2017} 
Ministerio de Derechos Sociales y Agenda 2030 y Fundación ONCE (2017) 
\textit{Estudio de accesibilidad universal en espacios públicos urbanizados y en la edificación en España}. 
Madrid: Observatorio de la Accesibilidad. 
Disponible en: \url{https://observatoriodelaaccesibilidad.es/wp-content/uploads/2024/04/Estudio-de-accesibilidad-universal-en-espacios-publicos-urbanizados-y-en-la-edificacion-en-espana-2017.pdf}.

\bibitem{rd_505_2007} 
Ministerio de la Presidencia (2007) 
«Real Decreto 505/2007, de 20 de abril, por el que se aprueban las condiciones básicas de accesibilidad y no discriminación de las personas con discapacidad para el acceso y utilización de los espacios públicos urbanizados y edificaciones», 
\textit{Boletín Oficial del Estado}, nº 113, págs. 20380--20387.

\bibitem{mitma_orden_tma851_2021} 
Ministerio de Transportes, Movilidad y Agenda Urbana (2021) 
«Orden TMA/851/2021, de 23 de julio, por la que se desarrollan las condiciones básicas de accesibilidad y no discriminación para el acceso y utilización de los espacios públicos urbanizados», 
\textit{Boletín Oficial del Estado}, nº 187, págs. 96328--96383.

\bibitem{oates_2006} 
Oates, B.J. (2006) 
\textit{Researching Information Systems and Computing}. 
London: Sage Publications.

\bibitem{once_pavimento_2021} 
ONCE (2021) 
\textit{Pavimento tacto-visual en el entorno urbano}. 
Madrid: Grupo Social ONCE.

\bibitem{once_plataformas_2024} 
ONCE (2024) 
\textit{Tipos de Plataformas Únicas. Condiciones de accesibilidad}. 
Madrid: Grupo Social ONCE.

\bibitem{predif_2016} 
PREDIF (2016) 
\textit{Criterios e indicadores de accesibilidad universal en entornos urbanos y turísticos}. 
Madrid: Plataforma Representativa Estatal de Personas con Discapacidad Física.

\bibitem{sala_alonso_2012} 
Sala Mozos, E. y Alonso López, F. (2012) 
\textit{La Accesibilidad Universal en los municipios: Guía para una política integral de promoción y gestión}. 
Madrid: IMSERSO - Ministerio de Sanidad, Servicios Sociales e Igualdad.

\bibitem{segittur_2021} 
SEGITTUR (2021) 
\textit{Turismo para todos: un reto para el sector y un derecho para las personas}. 
Madrid: Sociedad Mercantil Estatal para la Gestión de la Innovación y las Tecnologías Turísticas. 
Disponible en: \url{https://www.segittur.es/blog/sin-categorizar/turismo-para-todos-un-reto-para-el-sector-y-un-derecho-para-las-personas/}.

\bibitem{vazquez_ortega_2020} 
Vázquez Ortega, M. (2020) 
\textit{Accesibilidad en el 112: Diseño de un asistente virtual para personas sordas}. 
Trabajo Fin de Máster. Barcelona: Universitat Oberta de Catalunya (UOC).

\end{thebibliography}